%\documentclass[12pt]{article}
\documentclass[aps,prd,12pt,twocolumn,superscriptaddress,showkeys,floatfix]{revtex4-2}

\usepackage{kotex}
\usepackage{mathptmx}
\usepackage{amsmath}
\usepackage{booktabs}
\usepackage{graphicx}
\usepackage{setspace}
\usepackage[a4paper, margin=0.6in]{geometry}
\usepackage[hidelinks]{hyperref}

\setstretch{1.45}
\graphicspath{{artifacts/figures/predecoder_v2/}{artifacts/figures/predecoder/}{figure/}}

\newcommand{\predecoderfigure}[2]{%
\IfFileExists{artifacts/figures/predecoder_v2/#1_v2.pdf}{%
    \includegraphics[width=#2]{artifacts/figures/predecoder_v2/#1_v2.pdf}%
}{%
    \IfFileExists{artifacts/figures/predecoder_v2/#1_v2.png}{%
        \includegraphics[width=#2]{artifacts/figures/predecoder_v2/#1_v2.png}%
    }{%
        \IfFileExists{artifacts/figures/predecoder/#1.pdf}{%
            \includegraphics[width=#2]{artifacts/figures/predecoder/#1.pdf}%
        }{%
            \IfFileExists{artifacts/figures/predecoder/#1.png}{%
                \includegraphics[width=#2]{artifacts/figures/predecoder/#1.png}%
            }{%
                \fbox{\parbox{0.92\linewidth}{%
                \centering Figure file missing:\\
                \texttt{artifacts/figures/predecoder_v2/#1_v2.png}}}%
            }%
        }%
    }%
}%
}

\begin{document}

\title{Design and Evaluation of a Transition-Aware Neural Pre-Decoder for Surface-Code Quantum Error Correction}
\author{Sooho Chae}
\noaffiliation
\date{May 2026}

\begin{abstract}
표면 코드 양자 오류 정정에서 디코더는 측정된 syndrome으로부터 logical
frame을 추정한다. PyMatching은 이 문제에서 강한 기준 성능을 보이는
디코더이지만, 일부 샘플에서는 matching에 들어가기 전의 국소적인 syndrome
보정이 최종 판단을 개선할 수 있다. 이 논문은 PyMatching을 대체하지 않고,
그 앞단에서 작은 local detector edit을 선택하는 전이 정보 기반 신경망 사전
디코더를 설계하고 평가한다.

제안 모델은 36채널 syndrome/noise volume을 입력으로 받으며, 3차원 residual
trunk, local motif candidate 생성기, patch-head CandidateEditSelector,
benefit/harm 기반 ranking, 그리고 candidate-first selected-mode safety
policy로 구성된다. 최종 logical frame은 항상 PyMatching이 예측하며,
신경망은 수정된 syndrome을 전달할지 원본 syndrome을 유지할지만 결정한다.

Held-out correlated-noise family인 \texttt{stage\_c\_corr}에서 d3와 d5는 raw
PyMatching 대비 각각 \(+0.006591797\), \(+0.005615234\)의 selected-mode 평균
정확도 향상을 보였다. d3는 seed \(0..7\) 전체에서 양의 향상을 보였고, d5는
두 seed에서만 local edit을 채택하고 나머지 seed에서는 raw no-edit으로
fallback하는 보수적인 개선 양상을 보였다. 반면 d7의 향상은
\(+0.000151536\)에 머물렀다. d7의 모든 58개 seed에는 positive oracle
headroom이 남아 있었으므로, 실패 원인은 candidate 부족이 아니라 learned
selector의 ranking 및 generalization 한계로 해석된다.
\end{abstract}

\keywords{quantum error correction, surface code, neural pre-decoder, PyMatching, syndrome decoding, selector ranking}

\maketitle

\section{서론}

양자 상태는 계산 과정에서 gate error, measurement error, decoherence의 영향을
받는다. 작은 물리 오류도 누적되면 logical state가 손상되기 때문에, 실용적인
양자 계산에는 오류 정정이 필수적이다. 표면 코드는 국소 stabilizer 측정과
높은 오류 허용 특성 때문에 가장 널리 연구되는 오류 정정 코드 중 하나이다
~\cite{surfacecode}.

표면 코드에서 반복 측정으로 얻는 detector event는 오류가 발생했을 가능성이
있는 시공간적 위치 정보를 담고 있다. 디코더는 이 정보를 이용해 logical
effect를 추정한다. PyMatching은 minimum-weight perfect matching에 기반한
강력한 디코더이며~\cite{pymatching}, 본 논문에서도 최종 logical frame 추정은
PyMatching에 맡긴다.

이 연구의 출발점은 PyMatching 자체를 신경망으로 대체하는 것이 아니다.
대신 PyMatching이 받는 syndrome을 아주 작은 범위에서 미리 고칠 수 있는지
확인한다. 즉, 신경망은 ``어떤 logical class인가''를 직접 예측하지 않고,
``이 local detector edit을 적용하면 PyMatching의 판단이 나아지는가''를
판단한다. 이 관점은 direct neural decoder보다 보수적이지만, matching decoder의
안정성을 그대로 활용할 수 있다는 장점이 있다.

본 논문이 다루는 질문은 세 가지이다. 첫째, local syndrome edit이 실제로
PyMatching의 held-out 성능을 개선하는가. 둘째, 이 효과가 d3, d5, d7에서
어떻게 달라지는가. 셋째, d7에서 향상이 작다면 그 원인이 후보 공간의 부족인지,
아니면 후보를 고르는 selector의 일반화 문제인지 구분할 수 있는가.

\section{배경 및 관련 연구}

\subsection{표면 코드와 matching decoder}

표면 코드는 2차원 격자 위의 qubit에 대해 stabilizer 측정을 반복한다. 시간에
따라 변화하는 syndrome은 detector event로 표현되며, 디코더는 이 event pattern을
보고 correction의 logical effect를 추정한다. PyMatching은 이 문제를 graph
matching 문제로 변환하여 오류 chain을 찾는다. 해석 가능하고 구현이 안정적이기
때문에, 본 연구에서는 PyMatching을 기준 decoder이자 최종 decoder로 둔다.

\subsection{신경망 디코딩과 사전 디코딩}

신경망을 이용한 표면 코드 디코딩에는 크게 두 방향이 있다. 하나는 syndrome에서
logical class 또는 correction field를 직접 예측하는 방식이고, 다른 하나는
기존 디코더 앞에서 입력을 보정하는 방식이다. AlphaQubit은 direct neural
decoder의 대표적인 예이고~\cite{alphaqubit}, NVIDIA Ising-Decoding은
pre-decoder와 downstream decoder를 결합하는 구조적 선행 연구이다
~\cite{nvidiaising}. 본 프로젝트에서 시도한 direct neural baseline은 거리
증가에 따라 안정적인 성능을 보이지 못했다. 그 결과 최종 decoder를
교체하기보다, PyMatching 앞단의 보조 모듈로 신경망을 사용하는 방향이 더 방어
가능한 선택으로 남았다.

관련 연구와의 차이는 Table~\ref{tab:related-work}에 요약했다. 이 표는 동일
dataset에서의 직접 성능 비교가 아니라, 각 연구가 어떤 decoding 구조를 택했는지
구분하기 위한 문맥 비교이다. 따라서 본 논문의 수치는 본 논문에서 구성한
Stage A/B/C noise 환경과 held-out \texttt{stage\_c\_corr} 평가 안에서 해석해야
한다.

\begin{table*}[t]
\caption{관련 연구와 본 논문의 위치.}
\label{tab:related-work}
\centering
\small
\begin{tabular}{llll}
\toprule
work & decoder form & scope & relation \\
\midrule
AlphaQubit~\cite{alphaqubit} & direct neural decoder & large samples & replacement-style neural decoding \\
NVIDIA Ising-Decoding~\cite{nvidiaising} & AI pre-decoder & fast residual decoding & closest structural reference \\
Near-term neural decoder~\cite{neartermnn} & neural decoder & small experiments & near-term context \\
Google Willow~\cite{willowthreshold} & hardware decoding & below-threshold memory & hardware-scale context \\
This work & pre-decoder + PyMatching & d3/d5 gain, d7 limit & candidate/selected/oracle analysis \\
\bottomrule
\end{tabular}
\end{table*}

\section{데이터셋 및 노이즈 환경}

실험은 Stim 기반으로 직접 생성한 표면 코드 회로와 noise family를 바탕으로
수행했다~\cite{stim}. 주요 family는 \texttt{stage\_a\_si1000},
\texttt{stage\_b\_local}, \texttt{stage\_c\_corr}이다. 앞의 두 family는 학습과
검증에 사용했고, \texttt{stage\_c\_corr}는 최종 held-out 평가에 사용했다.

입력은 detector event만으로 구성하지 않았다. detector mask, geometry metadata,
noise-family indicator, distance/round statistics, event fraction,
physical-noise summary 등을 포함하여 36채널 syndrome/noise volume을 만들었다.
d3와 d5 성공 실험에서 tensor shape은 각각 \([36,4,4,4]\),
\([36,6,6,6]\)이다. 채널 정의는 거리와 무관하게 유지하고, spatial/temporal
volume 크기만 distance에 따라 달라진다.

\begin{table}[t]
\caption{실험에 사용한 noise family.}
\label{tab:noise-family}
\centering
\begin{tabular}{lll}
\toprule
family & role & use \\
\midrule
\texttt{stage\_a\_si1000} & baseline physical noise & train/validation \\
\texttt{stage\_b\_local} & local structured noise & train/validation \\
\texttt{stage\_c\_corr} & correlated noise & held-out evaluation \\
\bottomrule
\end{tabular}
\end{table}

\section{제안 방법}

제안하는 decoder는 전이 정보 기반 neural pre-decoder이다. 전체 흐름은
Figure~\ref{fig:pipeline}과 같다. 신경망은 local edit 후보를 평가하고, 선택된
edit을 syndrome에 적용할지 또는 raw no-edit을 유지할지 결정한다. 이후의
logical frame 추정은 PyMatching이 수행한다.

\begin{figure*}[t]
\centering
\predecoderfigure{fig1_predecoder_pipeline}{\textwidth}
\caption{제안하는 neural pre-decoder와 PyMatching 결합 구조. 신경망은 local
syndrome edit의 채택 여부만 결정하고, 최종 logical frame은 PyMatching이
예측한다.}
\label{fig:pipeline}
\end{figure*}

\subsection{SyndromeEditPreDecoder}

\texttt{SyndromeEditPreDecoder}는 3D convolution stem과 residual 3D convolution
block으로 구성된다. 성공한 d3/d5 recipe에서는 hidden channel \(24\), dense
hidden dimension \(64\), dropout \(0.1\)을 사용했다. trunk는 detector-level edit
logits, shot-level needs-edit logits, pooled shot feature를 만든다. 이 중 pooled
feature와 candidate feature가 patch-head selector의 입력으로 들어간다.

\subsection{Local motif candidate}

모델은 임의의 큰 syndrome edit을 허용하지 않는다. 성공 recipe는
\texttt{local\_motif\_selector}를 사용하며, local motif class와 top-k candidate
수를 제한한다. Identity, 즉 no-edit candidate는 항상 포함한다. 이 설계 덕분에
raw PyMatching은 외부 baseline일 뿐 아니라 selected-mode 내부의 안전 경로가
된다.

\subsection{Benefit/harm selector}

\texttt{CandidateEditSelector}는 각 local edit 후보에 score를 부여한다. 후보의
좋고 나쁨은 detector mask 복원 여부가 아니라, 그 edit을 적용한 뒤 PyMatching의
\texttt{logical\_class4} 판단이 raw no-edit보다 좋아지는지로 정의한다. 따라서
selector는 단순한 syndrome 복원기가 아니라, PyMatching의 최종 행동을 반영하는
decoder-aware ranking 모듈이다. Figure~\ref{fig:architecture}는 이 구조를
요약한다.

\begin{figure*}[t]
\centering
\predecoderfigure{fig2_model_architecture}{\textwidth}
\caption{\texttt{SyndromeEditPreDecoder}와 patch-head
\texttt{CandidateEditSelector}. Trunk가 edit/shot feature를 만들고, selector가
local candidate를 benefit/harm 관점에서 ranking한다.}
\label{fig:architecture}
\end{figure*}

\subsection{Selected-mode safety}

Candidate branch가 항상 최종 결과에 반영되는 것은 아니다. Validation evidence,
harm guard, support guard, plateau guard를 통과하지 못하면 selected mode는 raw
no-edit으로 fallback한다. 이 candidate-first safety policy는 성능을 높이기
위한 부가 장치가 아니라, harmful edit을 막기 위한 핵심 설계이다.

\section{실험 설정}

비교 대상은 네 가지이다. Raw PyMatching은 원본 syndrome을 그대로 사용한
결과이다. Selected predecoder는 safety policy가 선택한 edit 또는 no-edit을
적용한 결과이다. Candidate branch는 learned candidate를 fallback 없이 적용한
경우이고, target local-edit oracle은 후보 공간 안에서 가능한 최선의 local
edit 성능이다.

주요 metric은 held-out \texttt{stage\_c\_corr}에서의
\texttt{logical\_class4} accuracy이다. 또한 selected branch와 candidate branch를
분리하여, 성능 변화가 실제 채택된 안전한 edit에서 오는지, 아니면 후보 자체의
잠재력에서 오는지 구분했다.

\section{실험 결과}

\subsection{Main accuracy}

Table~\ref{tab:main-results}와 Figure~\ref{fig:main-accuracy}는 d3, d5, d7의
주요 결과를 보여준다. d3와 d5에서는 selected predecoder가 raw PyMatching보다
높은 정확도를 보였다. d7에서는 selected-mode 향상이 거의 0에 가까웠다.

\begin{table*}[t]
\caption{Held-out \texttt{stage\_c\_corr}에서의 main accuracy.}
\label{tab:main-results}
\centering
\begin{tabular}{lcccccc}
\toprule
distance & seeds & raw & selected & candidate & oracle & selected delta \\
\midrule
d3 & \(0..7\) & 0.928710938 & 0.935302734 & 0.935302734 & 0.992187500 & \(+0.006591797\) \\
d5 & \(0..7\) & 0.888671875 & 0.894287109 & 0.891845703 & 0.978515625 & \(+0.005615234\) \\
d7 & \(0..57\) & 0.873046875 & 0.873198411 & 0.871531519 & 0.984375000 & \(+0.000151536\) \\
\bottomrule
\end{tabular}
\end{table*}

\begin{figure*}[t]
\centering
\predecoderfigure{fig3_main_accuracy_comparison}{\textwidth}
\caption{거리별 held-out accuracy. d3/d5에서는 selected predecoder가 raw
PyMatching을 넘지만, d7에서는 oracle headroom이 큰데도 selected recovery가
거의 나타나지 않는다.}
\label{fig:main-accuracy}
\end{figure*}

\subsection{d3/d5 robustness}

d3는 seed \(0..7\) 전부에서 양의 selected delta를 보였다. d5는 두 seed에서만
local selector를 채택하고, 나머지 여섯 seed에서는 raw no-edit으로 fallback했다.
즉 d5의 평균 향상은 강한 전면적 개선이 아니라, 유익한 경우에만 edit을 채택한
보수적 개선이다.

\begin{table}[t]
\caption{d3/d5 seed \(0..7\) robustness summary.}
\label{tab:d3d5-robustness}
\centering
\begin{tabular}{lccc}
\toprule
distance & seed classes & mean delta & improved/harmed \\
\midrule
d3 & \(8/0/0\) & \(+0.006591797\) & \(205/151\) \\
d5 & \(2/6/0\) & \(+0.005615234\) & \(56/10\) \\
\bottomrule
\end{tabular}
\end{table}

Seed-level bootstrap 95\% confidence interval은 d3가
\([+0.004516602,+0.008544922]\), d5가
\([+0.000000000,+0.013671875]\)이다. Exact sign test에서도 d3의 one-sided
p-value는 \(0.003906250\)인 반면, d5는 \(0.250000000\)이다. 따라서 d3는
uniformly positive result로 제시할 수 있지만, d5는 positive mean과 selected
safety를 중심으로 해석해야 한다.

\begin{table}[t]
\caption{d3/d5 paired seed-level statistics.}
\label{tab:paired-stats}
\centering
\begin{tabular}{lccc}
\toprule
distance & seed classes & sign p & sign-flip p \\
\midrule
d3 & \(8/0/0\) & 0.003906250 & 0.003906250 \\
d5 & \(2/6/0\) & 0.250000000 & 0.250000000 \\
\bottomrule
\end{tabular}
\end{table}

\begin{table*}[t]
\caption{d5 seed-level selected-mode fallback behavior.}
\label{tab:d5-fallback}
\centering
\small
\begin{tabular}{clrrl}
\toprule
seed & selected mode & selected delta & candidate delta & interpretation \\
\midrule
0 & raw no-edit & \(+0.000000000\) & \(+0.000000000\) & fallback \\
1 & raw no-edit & \(+0.000000000\) & \(+0.000000000\) & fallback \\
2 & local selector & \(+0.021484375\) & \(+0.021484375\) & useful local edit \\
3 & local selector & \(+0.023437500\) & \(+0.023437500\) & useful local edit \\
4 & raw no-edit & \(+0.000000000\) & \(-0.007812500\) & harmful candidate blocked \\
5 & raw no-edit & \(+0.000000000\) & \(+0.000000000\) & fallback \\
6 & raw no-edit & \(+0.000000000\) & \(-0.011718750\) & harmful candidate blocked \\
7 & raw no-edit & \(+0.000000000\) & \(+0.000000000\) & fallback \\
\bottomrule
\end{tabular}
\end{table*}

Table~\ref{tab:d5-fallback}는 d5의 평균 향상이 seed 2와 seed 3에서 발생했음을
보여준다. 반대로 seed 4와 seed 6의 candidate branch는 harmful이었지만,
selected mode가 이를 raw no-edit으로 막았다.

\subsection{Candidate branch와 selected branch}

Candidate branch는 후보 공간의 잠재력을 보여주지만, fallback 없이 적용하면
harmful edit도 함께 들어간다. 특히 d7에서는 candidate branch의 improved/harmed
shot 수가 \(161/251\)로 harmful 쪽이 더 많다. Selected mode는 대부분 no-edit을
선택하여 큰 성능 하락을 막는다. 이 차이가 본 논문에서 candidate, selected,
oracle을 분리해 보고하는 이유이다.

\section{d7 한계 분석}

d7의 selected-mode delta는 \(+0.000151536\)에 불과하다. 그러나 target
local-edit oracle accuracy는 \(0.984375000\)이고, 모든 58개 seed에 positive
oracle headroom이 존재한다. 따라서 d7의 한계는 후보 공간 안에 유익한 local
edit이 없어서 생긴 문제가 아니다.

문제는 selector가 held-out 환경에서 유익한 후보를 안정적으로 고르지 못한다는
점이다. d7 candidate outcome은 positive \(6\), neutral \(35\), harmful \(17\)로
나뉘며, mean actual candidate delta는 \(-0.001515356\)이다.
Figure~\ref{fig:d7-oracle}은 oracle headroom과 false-positive 구조를 함께
보여준다.

\begin{figure*}[t]
\centering
\predecoderfigure{fig4_d7_oracle_gap_false_positive}{\textwidth}
\caption{d7의 oracle headroom과 false-positive bottleneck. 후보 공간에는 유익한
edit이 남아 있지만, learned selector가 validation-positive branch를 held-out에서
안정적으로 고르지 못한다.}
\label{fig:d7-oracle}
\end{figure*}

Validation-to-held-out mismatch는 이 문제를 더 분명하게 드러낸다.
Validation-positive candidate seed 22개 중 held-out에서 harmful이 된 seed는
13개이고, positive로 남은 seed는 5개뿐이다. False-positive ratio는
\(13/22=59.09\%\)이다.

\begin{table}[t]
\caption{d7 validation class와 held-out outcome의 불일치.}
\label{tab:d7-mismatch}
\centering
\begin{tabular}{lccc}
\toprule
validation class & harmful & neutral & positive \\
\midrule
neutral & 4 & 31 & 1 \\
positive & 13 & 4 & 5 \\
\bottomrule
\end{tabular}
\end{table}

\begin{figure*}[t]
\centering
\predecoderfigure{fig5_d7_validation_heldout_scatter}{\textwidth}
\caption{d7 validation delta와 held-out candidate delta의 seed-level scatter.
Validation-positive evidence가 held-out correlated noise에서 그대로 유지되지
않는다는 점이 핵심이다.}
\label{fig:d7-scatter}
\end{figure*}

Harmful candidate seed는 총 17개였고, selected-mode guard는 이 \(17/17\)개를
모두 raw no-edit으로 fallback시켰다. 즉 selected-mode safety는 d7에서 성능을
크게 올리지는 못했지만, harmful candidate를 실제 출력으로 내보내는 것은 막았다.

추가적인 d7 threshold sweep도 검토했지만, 현재 model family 안에서는 합리적인
해결책을 찾지 못했다. Scalar adoption-threshold grid에서는 \(183040\)개 policy
중 preserve/recover/block sentinel gate를 통과한 policy가 없었다. Cross-family
hard positive-vs-negative objective와 candidate-compatibility pairwise top-k도
true positive 보존과 false positive 차단을 동시에 만족하지 못했다. 따라서 d7의
다음 개선은 단순 threshold 조정이 아니라 selector ranking/generalization
목표를 다시 설계하는 방향이어야 한다.

\begin{figure*}[t]
\centering
\predecoderfigure{fig6_oracle_recovery_distribution}{\textwidth}
\caption{Seed-level oracle-gap recovery. d3/d5는 일부 oracle gap을 회복하지만,
d7 selected mode는 candidate-oracle headroom을 거의 활용하지 못한다.}
\label{fig:oracle-recovery}
\end{figure*}

\section{논의}

이 논문의 핵심 결과는 두 부분으로 나뉜다. 첫째, d3와 d5에서는 neural
pre-decoder가 PyMatching의 front-end로 실제 이득을 줄 수 있다. 둘째, d7에서는
local edit headroom이 남아 있어도 selector가 이를 안정적으로 선택하지 못하면
성능 향상으로 이어지지 않는다.

이 구분은 중요하다. 단순히 ``d7에서 실패했다''고만 쓰면 후보 공간 자체가
부족한 것처럼 보일 수 있다. 그러나 oracle 분석은 그렇지 않음을 보여준다.
후보 공간에는 유익한 edit이 있고, 병목은 그 후보를 held-out noise에서
구분하는 ranking 문제에 있다. 따라서 본 연구의 d7 결과는 부정적 결과라기보다,
현재 모델을 더 큰 거리로 확장할 때 먼저 해결해야 할 지점을 보여주는 분석
결과로 보는 것이 타당하다.

한계도 분명하다. d7에서 robust learned recovery는 달성하지 못했다. 또한
결과는 본 논문에서 구성한 noise family와 seed 범위 안에서 해석해야 한다.
\texttt{config.py}에는 Stage D leakage surrogate와 Stage E DQLR 설정 스키마가
있지만, 이를 shot-level dynamics로 적용하는 \texttt{postprocess.py} 모듈은
아직 없다. \texttt{noise\_willowcore.py}도 Stage D/E 요청을 명시적으로
차단한다. 따라서 본 논문의 정량 평가는 Stage A/B/C 범위로 제한하며, leakage와
DQLR까지 포함하는 평가는 후속 연구로 남긴다.

\section{결론}

이 논문은 PyMatching 앞단에서 local syndrome edit을 선택하는 전이 정보 기반
신경망 사전 디코더를 설계하고 평가했다. 제안한 구조는 PyMatching을 대체하지
않고, selected local edit 또는 raw no-edit syndrome을 PyMatching에 넘기는
방식으로 동작한다.

Held-out \texttt{stage\_c\_corr} 평가에서 d3와 d5는 raw PyMatching 대비 각각
\(+0.006591797\), \(+0.005615234\)의 selected-mode 평균 향상을 보였다. d3는
검사한 모든 seed에서 양의 결과를 보였고, d5는 두 seed의 local edit 채택과
나머지 seed의 fallback으로 구성된 보수적 개선이었다.

d7에서는 selected-mode 향상이 거의 없었다. 그러나 oracle과 candidate-oracle
분석은 유익한 local edit 후보가 여전히 존재함을 보였다. 따라서 현재 한계는
candidate coverage가 아니라 selector ranking과 generalization에 있다. 이
결론은 neural pre-decoder를 더 큰 거리로 확장하려면 단순 threshold tuning보다
후보 ranking objective의 재설계가 필요하다는 점을 시사한다.

\section*{Reproducibility Note}

결과표는 \texttt{PREDECODER\_FINAL\_RESULT\_TABLES.md},
\texttt{PREDECODER\_D3\_D5\_PAIRED\_STATISTICS.md},
\texttt{PREDECODER\_ORACLE\_RECOVERY\_DISTRIBUTION.md},
\texttt{PREDECODER\_D7\_HARMFUL\_EDIT\_TAXONOMY.md},
\texttt{PREDECODER\_HYPERPARAMETER\_SENSITIVITY.md}, 그리고
\texttt{PREDECODER\_REPRODUCIBILITY\_PACKAGE.md}를 기준으로 한다. 최종
consistency check는
\texttt{artifacts/eval/nn/sedp\_final\_result\_consistency\_check.json}에
\(37/37\) checks, \(0\) failures로 기록되어 있다.

\bibliographystyle{apsrev4-2}
\bibliography{main}

\end{document}
